% Original teaching example; the tree is supplied, not claimed to be greedy CART.
% Fixed-tree, unweighted squared loss: leaf predictions are response means.
% Canvas: 158 mm x 64 mm. Plot window: 0 <= x <= 6; both variables dimensionless.
% Dependencies: project palette and TikZ with arrows.meta; no external resources.
\begin{tikzpicture}[x=1mm,y=1mm,font=\footnotesize,text=FigureOutput,
  tree rule/.style={draw=FigureModel,fill=FigureModel!8,
    line width=0.6pt,minimum width=18mm,minimum height=7mm,
    inner sep=1mm,align=center},
  tree leaf/.style={draw=FigureOutput!55,fill=FigureOutput!4,
    line width=0.5pt,minimum width=23mm,minimum height=10mm,
    inner sep=1mm,align=center},
  tree edge/.style={-{Stealth[length=1.6mm]},draw=FigureMuted,
    line width=0.6pt},
  prediction/.style={draw=FigureOutput,line width=1pt,
    dash pattern=on 2.4pt off 1.4pt},
  split guide/.style={draw=FigureModel,line width=0.6pt},
  plot axis/.style={-{Stealth[length=1.5mm]},draw=FigureMuted,
    line width=0.5pt}]
  \path[use as bounding box] (0,0) rectangle (158,64);

  % Samples are ordered by x. Keep the leaf sample labels consistent when editing.
  \def\TreeSamples{0.5/1,1.5/2,2.5/4,3.5/5,4.5/2,5.5/3}
  \def\TreeThresholdOne{2}
  \def\TreeThresholdTwo{4}
  \pgfmathsetmacro{\TreeMeanOne}{(1+2)/2}
  \pgfmathsetmacro{\TreeMeanTwo}{(4+5)/2}
  \pgfmathsetmacro{\TreeMeanThree}{(2+3)/2}

  \node[anchor=west,inner sep=0pt] at (1,61) {(a) 给定阈值树};
  \node[anchor=west,inner sep=0pt] at (82,61) {(b) 一维分段常数预测};
  \draw[FigureGrid,line width=0.5pt] (75,6)--(75,58);
  \node[text=FigureMuted,inner sep=0pt] at (37.5,55.5)
    {平方损失：固定树上的叶内均值};

  \node[tree rule] (root) at (33,48.5) {$x\leq\TreeThresholdOne$？};
  \node[tree leaf] (leaf1) at (13,30)
    {$R_1:\ c_1=\pgfmathprintnumber{\TreeMeanOne}$\\
     \textcolor{FigureInput}{$y_1=1,\ y_2=2$}};
  \node[tree rule] (split2) at (49,30) {$x\leq\TreeThresholdTwo$？};
  \node[tree leaf] (leaf2) at (36.5,12)
    {$R_2:\ c_2=\pgfmathprintnumber{\TreeMeanTwo}$\\
     \textcolor{FigureInput}{$y_3=4,\ y_4=5$}};
  \node[tree leaf] (leaf3) at (61.5,12)
    {$R_3:\ c_3=\pgfmathprintnumber{\TreeMeanThree}$\\
     \textcolor{FigureInput}{$y_5=2,\ y_6=3$}};

  \draw[tree edge] (root.south west)
    --node[pos=0.5,left,inner sep=1pt] {是} (leaf1.north);
  \draw[tree edge] (root.south east)
    --node[pos=0.5,right,inner sep=1pt] {否} (split2.north);
  \draw[tree edge] (split2.south west)
    --node[pos=0.5,left,inner sep=1pt] {是} (leaf2.north);
  \draw[tree edge] (split2.south east)
    --node[pos=0.5,right,inner sep=1pt] {否} (leaf3.north);
  \node[text=FigureMuted,inner sep=0pt] at (37.5,2.8)
    {等于阈值时走“是”分支};

  \fill[FigureInput] (83,56) circle[radius=0.75mm];
  \node[anchor=west,inner sep=0pt] at (86,56) {样本$(x_i,y_i)$};
  \draw[prediction] (119,56)--(127,56);
  \node[anchor=west,inner sep=0pt] at (130,56) {预测$f_T(x)$};
  \node[anchor=west,inner sep=0pt] at (82,51.7)
    {$y,\ f_T(x)$（无量纲）};
  \node[inner sep=0pt] at (120,2.8) {$x$（无量纲）};

  % One x unit is 12 mm; one y unit is 7 mm. Coordinates below are data values.
  \begin{scope}[shift={(84,12)},x=12mm,y=7mm]
    \foreach \tick in {1,2,3,4,5}{
      \draw[FigureGrid,line width=.5pt] (0,\tick)--(6,\tick);
    }
    \foreach \threshold in {\TreeThresholdOne,\TreeThresholdTwo}{
      \draw[split guide] (\threshold,-0.1)--(\threshold,5.3);
      \node[below=0.7mm,inner sep=0pt,text=FigureModel]
        at (\threshold,-0.1) {$\threshold$};
    }
    \draw[plot axis] (0,0)--(6.03,0);
    \draw[plot axis] (0,0)--(0,5.4);
    \foreach \tick in {0,1,2,3,4,5}{
      \draw[FigureMuted,line width=.5pt] (0,\tick)--(-0.06,\tick)
        node[left=0.8mm,inner sep=0pt] {$\tick$};
    }
    \foreach \samplex/\sampley in \TreeSamples {
      \draw[FigureMuted,line width=.5pt] (\samplex,0)--(\samplex,-0.1)
        node[below=0.7mm,inner sep=0pt,text=FigureMuted] {$\samplex$};
    }

    % Separate segments avoid suggesting that vertical jumps are function values.
    \draw[prediction] (0,\TreeMeanOne)--(\TreeThresholdOne,\TreeMeanOne);
    \draw[prediction] (\TreeThresholdOne,\TreeMeanTwo)
      --(\TreeThresholdTwo,\TreeMeanTwo);
    \draw[prediction] (\TreeThresholdTwo,\TreeMeanThree)--(6,\TreeMeanThree);

    \foreach \samplex/\sampley in \TreeSamples {
      \fill[FigureInput] (\samplex,\sampley) circle[radius=0.75mm];
    }
    % Equality belongs to the interval on the left: closed above/below as needed.
    \fill[FigureOutput] (\TreeThresholdOne,\TreeMeanOne) circle[radius=0.8mm];
    \fill[FigureOutput] (\TreeThresholdTwo,\TreeMeanTwo) circle[radius=0.8mm];
    \draw[draw=FigureOutput,fill=FigureOutput!0,line width=0.7pt]
      (\TreeThresholdOne,\TreeMeanTwo) circle[radius=0.8mm];
    \draw[draw=FigureOutput,fill=FigureOutput!0,line width=0.7pt]
      (\TreeThresholdTwo,\TreeMeanThree) circle[radius=0.8mm];

    \node[inner sep=0pt] at (1,0.45)
      {$R_1:\ c_1=\pgfmathprintnumber{\TreeMeanOne}$};
    \node[inner sep=0pt] at (3,3.25)
      {$R_2:\ c_2=\pgfmathprintnumber{\TreeMeanTwo}$};
    \node[inner sep=0pt] at (5,1.2)
      {$R_3:\ c_3=\pgfmathprintnumber{\TreeMeanThree}$};
  \end{scope}
\end{tikzpicture}
